% Generated draft: requires mathematical and source review.
\documentclass[11pt]{article}
\usepackage[margin=25mm]{geometry}
\usepackage{fontspec}
\setmainfont{Latin Modern Roman}
\usepackage{amsmath,amssymb,mathtools,unicode-math}
\setmathfont{Latin Modern Math}
\usepackage{enumitem,needspace,xurl,hyperref}
\setlength{\parindent}{0pt}
\setlength{\parskip}{5pt}
\setlength{\emergencystretch}{3em}
\newcommand{\N}{\mathbb{N}}
\newcommand{\Z}{\mathbb{Z}}
\newcommand{\Q}{\mathbb{Q}}
\newcommand{\R}{\mathbb{R}}
\newcommand{\C}{\mathbb{C}}
\newcommand{\F}{\mathbb{F}}
\newcommand{\A}{\mathbb{A}}
\newcommand{\PP}{\mathbb{P}}
\newcommand{\E}{\mathbb{E}}
\newcommand{\eps}{\varepsilon}
\newcommand{\dd}{\,\mathrm{d}}
\begin{document}
\section*{UnsolvedMath ID: 30005648}
\textbf{Source problem code:} catalogue entry
\subsection{Definitions and mathematical statement}
Let \(K\) denote the field in the supplied record, \(X/K\) a geometrically connected curve of genus \(g>1\), \(J\) its Jacobian, and \(x_0\in X(K)\). The record uses the Albanese embedding
\[
\iota_{x_0}\colon X\hookrightarrow J,\qquad x\longmapsto [x-x_0].
\]
Define
\[
J(K)_{\mathrm{tors}}=\{P\in J(K): [n]P=0\text{ for some integer }n\geq 1\},
\qquad
T_K=\iota_{x_0}(X)(K)\cap J(K)_{\mathrm{tors}}.
\]
Here \([n]\) denotes multiplication by \(n\) on \(J\). Thus \(T_K\) is the source's \(X_{\mathrm{tors}}\), expressed as a subset of the Jacobian using the specified base point.

The supplied question is: for every such \(X,J,x_0\), if \(J\) has supersingular good reduction over \(K\), does
\[
\forall\sigma\in I_K\;\forall P\in T_K,\qquad \sigma(P)=P
\]
hold?

This formulation has unresolved qualifications. The record does not specify the field \(K\), a valuation, its residue characteristic, or the definition of \(I_K\). In the usual discretely valued setting, good reduction means that \(J\) extends to an abelian scheme over the valuation ring; supersingular good reduction means that its special fiber becomes isogenous over an algebraic closure of the residue field to a product of supersingular elliptic curves. An inertia subgroup is defined relative to the valuation and a choice of its prolongation to a separable closure. These explanations describe the conventional setting, but that setting is not supplied as an explicit hypothesis.

There is also a substantive point-domain ambiguity: the displayed set consists of \(K\)-rational points, which are fixed by the usual action of \(\operatorname{Gal}(\overline K/K)\). Consequently, the supplied rational-point notation does not specify the intended nontrivial inertia question. A possible intended set is
\[
T_{\overline K}=\iota_{x_0}(X)(\overline K)\cap J(\overline K)_{\mathrm{tors}},
\]
with the question whether inertia fixes every element of this set, but replacing \(T_K\) by \(T_{\overline K}\) is an unverified correction and is not adopted here. The source also does not explicitly state smoothness or properness hypotheses for \(X\). These issues require human mathematical review before a definitive formulation can be accepted.
\subsection{Short English statement}
The supplied problem asks whether supersingular good reduction of a curve's Jacobian forces inertia to fix all torsion points lying on the curve embedded using a chosen rational base point. The field and valuation hypotheses are missing, and the written set uses \(K\)-rational points, whose usual Galois action is already trivial. The intended geometric-point formulation therefore remains unresolved. The supplied dated curation labels the problem partially solved.
\subsection{Sources}
\begin{itemize}
\item Local UnsolvedMath registry, permanent ID \(30005648\), source problem code \texttt{catalogue entry}, title ``Inertia Actions on Torsion Points under Supersingular Reduction.'' The supplied record belongs to the Oberwolfach Reports Open Problems collection. No individual proposer is supplied.
\item Supplied source citation: MFO--RIMS Tandem Workshop: Arithmetic Homotopy and Galois Theory (2024), \emph{Oberwolfach Reports}, DOI \texttt{10.4171/owr/2023/42}, pp.~2377--2488. \url{https://doi.org/10.4171/owr/2023/42}. The record separately gives proposed year 2023; the relationship between these year labels is not resolved here.
\item Supplied literature reference: ``Takao, Coleman conjecture report.'' \url{https://doi.org/10.4171/owr/2023/42}. This supplied attribution does not identify an individual proposer of the problem.
\item Supplied literature reference: ``Hoshi, torsion-point ramification results.'' \url{https://www.kurims.kyoto-u.ac.jp/~yuichiro/papers_e.html}. The supplied link is a papers page rather than an identified article or theorem citation.
\item Literature assessment dated exactly \texttt{2026-08-22}: the supplied curation labels the problem partially solved and states that inertia triviality is proved in the superspecial case and for genus two with supersingular good reduction, while the general higher-genus supersingular-good-reduction question remains open. This is supplied dated curation derived from the OWR record, not an independently verified assessment; its status and cited claims require independent verification.
\item The registry labels the statement \texttt{corrected\_verified} and reports that it was checked against the official Oberwolfach report. These are supplied verification claims, not verification performed for this draft, and do not resolve the formulation issues recorded above.
\end{itemize}
\end{document}
